\chapter*{Article information and Abstract／论文信息与摘要}
\addcontentsline{toc}{chapter}{Article information and Abstract／论文信息与摘要}
\EN{Riemann--Hilbert Approach for the Hirota--Satsuma Coupled KdV System}
\CN{Hirota--Satsuma 耦合 KdV 系统的 Riemann--Hilbert 方法}
\EN{Haibing Zhang; Xianguo Geng; Huan Liu.}
\CN{作者：Haibing Zhang、Xianguo Geng、Huan Liu。}
\EN{School of Mathematics and Statistics, Zhengzhou University, Zhengzhou, Henan, China; Institute of Mathematics, Henan Academy of Sciences, Zhengzhou, Henan, China.}
\CN{作者单位：中国河南郑州，郑州大学数学与统计学院；中国河南郑州，河南省科学院数学研究所。}
\EN{Correspondence: Xianguo Geng (\texttt{xggeng@zzu.edu.cn}).}
\CN{通讯作者：Xianguo Geng，联系邮箱同上。}
\EN{Received: 29 April 2025. Revised: 22 August 2025. Accepted: 1 September 2025.}
\CN{收稿：2025 年 4 月 29 日；修订：2025 年 8 月 22 日；接受：2025 年 9 月 1 日。}
\EN{Funding: This work was supported by the National Natural Science Foundation of China (Grant Nos. 12471234, 12171439) and Science Foundation of Henan Academy of Sciences (Grant No. 20252319002).}
\CN{资助：国家自然科学基金（项目编号 12471234、12171439）与河南省科学院科学基金（项目编号 20252319002）。}
\EN{Keywords: Hirota--Satsuma coupled KdV system; Riemann--Hilbert problem; soliton solutions.}
\CN{关键词：Hirota--Satsuma 耦合 KdV 系统；Riemann--Hilbert 问题；孤子解。}
\EN{In this paper, the Riemann--Hilbert (RH) method is developed to solve the Hirota--Satsuma coupled KdV (HScKdV) system associated with $4\times4$ matrix spectral problem.}
\CN{本文发展 RH 方法，以求解与四阶矩阵谱问题相联系的 HScKdV 系统。}
\EN{Because the spectral matrix is asymmetric and high order, it brings great difficulties to analysis and solution.}
\CN{谱矩阵的不对称结构和较高阶数，使分析与求解面临较大困难。}
\EN{First, a direct scattering problem is carried out, from which the initial data are mapped to the scattering data.}
\CN{首先研究直接散射问题，将初始数据映射为散射数据。}
\EN{On the basis of introducing the generalized cross product of vectors, the basic meromorphic matrix eigenfunctions of corresponding Lax pairs are expressed by the Jost solutions and adjoint Jost solutions.}
\CN{引入向量的广义叉积后，用 Jost 解及伴随 Jost 解表达相应 Lax 对的基本亚纯矩阵特征函数。}
\EN{Then, the inverse scattering problem is characterized as the $4\times4$ matrix RH problem, which gives the formula for constructing solutions of the HScKdV system.}
\CN{继而将逆散射问题表述为四阶矩阵 RH 问题，由此获得构造 HScKdV 系统解的公式。}
\EN{As an example, the RH problem corresponding to the reflectionless case is solved and the soliton solution of the HScKdV system is obtained.}
\CN{作为实例，求解无反射情形的 RH 问题，并得到 HScKdV 系统的孤子解。}
\EN{Mathematics Subject Classification: 35Q15, 35B40.}
\CN{数学主题分类：35Q15、35B40。}
\begin{note}
摘要给出的是研究路线，具体断言的条件和证明见后文相应结果。
“直接散射”从位势求谱数据；“逆散射”由谱数据恢复位势。
“亚纯”允许在离散谱处有极点，与在全部谱平面上解析不同。
无反射时消失的是连续谱的反射数据，离散谱仍能产生孤子，故无反射不等于零位势。
本文研究范围限于无孤子与纯孤子两类，并未完成含连续谱和离散谱的一般混合情形，也未在本文给出长时间渐近定理。
\end{note}
